% TEMA 1: PILOTO DIDACTICO PARA PROYECCION EN CLASE
% Archivo independiente: no sustituye temas/tema1.tex.
% Compilar desde la raiz: latexmk -pdf -outdir=pdf temas/tema1_piloto_didactico.tex
% El estilo comun y el formato 16:9 (12 pt) se heredan de main.tex.
\documentclass[main.tex]{subfiles}

\begin{document}
% Marcadores editables para fotografias/ilustraciones aun no incorporadas.
% Uso: \pilotoimagen{anchura}{altura}{descripcion}
\newcommand{\pilotoimagen}[3]{%
  \begin{center}
  \fcolorbox{mateBlueDark}{mateBlueInner}{%
    \begin{minipage}[c][#2][c]{#1}
      \centering\color{mateBlueDark}\sffamily
      {\bfseries IMAGEN PENDIENTE}\par\smallskip
      {\footnotesize #3}
    \end{minipage}%
  }%
  \end{center}%
}

\graphicspath{{imagenes/tema1/}{../imagenes/tema1/}}

% Portada automatica del estilo original.
\portadatema{1}{Funciones de varias variables}

% ===== BLOQUE A. UN PROBLEMA QUE NECESITA DOS VARIABLES =====
% 1
\begin{frame}{¿Cómo representarías un terreno digital?}
\begin{columns}[c]
\begin{column}{.53\textwidth}
Un videojuego necesita conocer la \textbf{altura del suelo} en cada posición del mapa.
\medskip

Si conocemos las coordenadas $x$ e $y$, ¿podemos calcular la altura $z$?
\medskip

\alert{Pregunta para clase:} ¿cuáles son las entradas y cuál es la salida?
\end{column}
\begin{column}{.44\textwidth}
\pilotoimagen{.88\linewidth}{36mm}{Terreno 3D de un videojuego con ejes $x$, $y$, $z$}
\end{column}
\end{columns}
\end{frame}

% 2
\begin{frame}{De una variable a dos}
Ya conocemos funciones como $y=g(x)$: a cada valor de $x$ le asignamos un resultado.
\bigskip

Para un terreno, la altura depende de \textbf{dos coordenadas}:
\[
\boxed{z=h(x,y)}
\]
\medskip

$x$ e $y$ son las \textbf{variables independientes}; $z$ es la \textbf{variable dependiente}.
\end{frame}

% 3
\begin{frame}{Construyamos nuestra montaña}
Un modelo matemático sencillo de una colina es
\[
\boxed{h(x,y)=\sqrt{9-x^2-y^2}.}
\]
\medskip

¿Qué altura obtenemos en estos puntos?
\[
\begin{array}{c|c}
(x,y)&h(x,y)\\\hline
(0,0)&3\\
(1,2)&2\\
(3,2)&\text{¿?}
\end{array}
\]
\end{frame}

% 4
\begin{frame}{¿Qué es una función de dos variables?}
Una función $f$ asigna a cada par $(x,y)$ de su \textbf{dominio} un único número real:
\[
 f:D\subseteq\mathbb R^2\longrightarrow\mathbb R,
 \qquad (x,y)\longmapsto f(x,y).
\]
\medskip

\textbf{En nuestro ejemplo:} $(x,y)$ es una posición del mapa y $h(x,y)$ es su altura.
\medskip

No todos los pares $(x,y)$ tienen por qué ser entradas válidas.
\end{frame}

% 5
\begin{frame}{Pregunta rápida: ¿un valor o muchos?}
En la función $h(x,y)$:
\begin{enumerate}
\item ¿Una misma posición $(x,y)$ puede tener dos alturas distintas?
\medskip
\item ¿Dos posiciones diferentes pueden tener la misma altura?
\end{enumerate}
\vfill
\alert{Piensa 30 segundos y justifica las respuestas.}
\end{frame}

% ===== BLOQUE B. DOMINIO, IMAGEN Y CODOMINIO =====
% 6
\begin{frame}{¿Por qué falla una de las posiciones?}
En el punto $(3,2)$:
\[
 h(3,2)=\sqrt{9-3^2-2^2}=\sqrt{-4}.
\]
\medskip

No es un número real. Por tanto, nuestro modelo \textbf{no está definido} en esa posición.
\medskip

\textbf{Nueva pregunta:} ¿qué puntos del plano son válidos?
\end{frame}

% 7
\begin{frame}{El dominio de nuestra montaña}
Para calcular $h(x,y)=\sqrt{9-x^2-y^2}$ necesitamos
\[
9-x^2-y^2\geq0
\quad\Longleftrightarrow\quad
x^2+y^2\leq9.
\]
\[
\boxed{D=\{(x,y)\in\mathbb R^2:x^2+y^2\leq9\}.}
\]
\medskip

Es un \textbf{disco de radio 3}, incluida la circunferencia frontera.
\end{frame}

% 8
\begin{frame}{Un dominio se puede dibujar}
\begin{columns}[c]
\begin{column}{.49\textwidth}
\centering
\begin{tikzpicture}[scale=.72]
  \draw[->,gray] (-3.6,0)--(3.65,0) node[right] {$x$};
  \draw[->,gray] (0,-3.55)--(0,3.6) node[above] {$y$};
  \fill[mateBlueInner] (0,0) circle (3);
  \draw[very thick,mateBlueDark] (0,0) circle (3);
  \fill[mateBlueDark] (0,0) circle (.055);
  \node[below right] at (3,0) {$3$};
  \node[above left] at (0,3) {$3$};
  \node[below left] at (-3,0) {$-3$};
\end{tikzpicture}
\end{column}
\begin{column}{.49\textwidth}
\textbf{Lo que representa el dibujo:}
\begin{itemize}
\item Cada punto del disco es una entrada válida.
\item La frontera también pertenece al dominio.
\item Fuera del disco no hay altura real.
\end{itemize}
\end{column}
\end{columns}
\end{frame}

% 9
\begin{frame}{Dominio, codominio e imagen no son lo mismo}
Para $h(x,y)=\sqrt{9-x^2-y^2}$:
\begin{itemize}
\item \textbf{Dominio:} posiciones válidas, $D=\{x^2+y^2\leq9\}$.
\medskip
\item \textbf{Codominio:} podemos elegir $\mathbb R$ como conjunto de salida.
\medskip
\item \textbf{Imagen:} alturas que realmente aparecen, $[0,3]$.
\end{itemize}
\medskip

\alert{Comprueba:} ¿dónde está la altura máxima?, ¿y la mínima?
\end{frame}

% 10
\begin{frame}{Otro dominio: ahora hay un logaritmo}
Considera la función
\[
f(x,y)=\ln(y-x^2).
\]
\medskip

\textbf{Pista:} el argumento de un logaritmo real debe ser positivo.
\[
y-x^2>0
\quad\Longleftrightarrow\quad
\boxed{y>x^2.}
\]
\medskip

El dominio es la región \textbf{por encima} de la parábola $y=x^2$, sin incluir la curva.
\end{frame}

% 11
\begin{frame}{Minirreto: una división prohibida}
\[
 f(x,y)=\frac{3x-5y}{y-x^2}
\]
\medskip

\textbf{Antes de mirar la solución:} ¿qué curva debemos excluir?
\pause
\[
\boxed{D=\{(x,y)\in\mathbb R^2:y\ne x^2\}.}
\]
\medskip

No se puede dividir entre cero: excluimos los puntos de la parábola $y=x^2$.
\end{frame}

% ===== BLOQUE C. DE 2D A 3D =====
% 12
\begin{frame}{Del mapa plano al paisaje 3D}
Cada entrada $(x,y)$ produce una altura $z=h(x,y)$.
\[
(x,y)\longmapsto(x,y,h(x,y)).
\]
\begin{columns}[c]
\begin{column}{.50\textwidth}
\pilotoimagen{.89\linewidth}{32mm}{Mapa 2D con dominio circular y ejes $x,y$}
\end{column}
\begin{column}{.50\textwidth}
\pilotoimagen{.89\linewidth}{32mm}{Semiesfera superior: gráfica de $h(x,y)$ en 3D}
\end{column}
\end{columns}
\end{frame}

% 13
\begin{frame}{La gráfica no es el dominio}
Si $f:D\to\mathbb R$, su \textbf{gráfica} es
\[
G_f=\{(x,y,z)\in\mathbb R^3:(x,y)\in D,\ z=f(x,y)\}.
\]
\medskip

\begin{itemize}
\item El \textbf{dominio} está en el plano $XY$.
\item La \textbf{gráfica} suele ser una superficie en el espacio.
\end{itemize}
\medskip

Una misma posición del dominio corresponde a \textbf{un único punto} de esa gráfica.
\end{frame}

% 14
\begin{frame}{No todas las superficies son montañas}
\begin{columns}[c]
\begin{column}{.55\textwidth}
\[
f(x,y)=x+y^2.
\]
\medskip

Aquí el dominio natural es todo $\mathbb R^2$.
\medskip

\textbf{Pregunta:} ¿la altura es siempre positiva? ¿La superficie tiene una cima?
\end{column}
\begin{column}{.43\textwidth}
\pilotoimagen{.9\linewidth}{35mm}{Gráfica 3D de $z=x+y^2$ (imagen original disponible en el tema 1)}
\end{column}
\end{columns}
\end{frame}

% ===== BLOQUE D. CURVAS DE NIVEL =====
% 15
\begin{frame}{¿Y si queremos un mapa sin perspectiva 3D?}
Imagina que el agua alcanza la altura de \textbf{2 metros}.
\medskip

¿Qué puntos de nuestra montaña quedan exactamente a esa altura?
\[
\sqrt{9-x^2-y^2}=2.
\]
\vfill
\alert{Piensa: ¿será un punto, una recta o una circunferencia?}
\end{frame}

% 16
\begin{frame}{Resolvamos la pregunta anterior}
Buscamos todos los puntos donde $h(x,y)=2$:
\begin{align*}
\sqrt{9-x^2-y^2}&=2,\\
9-x^2-y^2&=4,\\
\boxed{x^2+y^2=5}
\end{align*}
\medskip

Son los puntos de una \textbf{circunferencia de radio $\sqrt5$}.
\end{frame}

% 17
\begin{frame}{Distintas alturas, distintas curvas}
\begin{columns}[c]
\begin{column}{.53\textwidth}
\centering
\begin{tikzpicture}[scale=.76]
  \draw[->,gray] (-3.65,0)--(3.8,0) node[right] {$x$};
  \draw[->,gray] (0,-3.6)--(0,3.75) node[above] {$y$};
  \draw[thick,mateBlueOuter] (0,0) circle (3);
  \draw[very thick,mateBlueDark] (0,0) circle (2.82843);
  \draw[very thick,cyan!65!black] (0,0) circle (2.23607);
  \fill[mateBlueDark] (0,0) circle (.1);
  \node[right,fill=white,inner sep=1pt,font=\scriptsize] at (3,0) {$k=0$};
  \node[above right,fill=white,inner sep=1pt,font=\scriptsize] at (2,2) {$k=1$};
  \node[above right,fill=white,inner sep=1pt,font=\scriptsize] at (1.4,1.74) {$k=2$};
  \node[below right,font=\scriptsize] at (.14,-.3) {$k=3$};
\end{tikzpicture}
\end{column}
\begin{column}{.45\textwidth}
Para $0\le k\le3$:
\[
x^2+y^2=9-k^2.
\]
\small
\begin{tabular}{c|c}
$k$ & radio \\\hline
$0$ & $3$\\
$1$ & $\sqrt8$\\
$2$ & $\sqrt5$\\
$3$ & $0$\\
\end{tabular}
\medskip

Al subir, las curvas se hacen más pequeñas.
\end{column}
\end{columns}
\end{frame}

% 18
\begin{frame}{Ahora sí: definición de curva de nivel}
Una \textbf{curva de nivel} de $f(x,y)$ es el conjunto de puntos de su dominio donde la función toma un mismo valor $k$:
\[
\boxed{\{(x,y)\in D:f(x,y)=k\}.}
\]
\medskip

Al cambiar $k$, obtenemos diferentes curvas (o conjuntos de nivel, que pueden ser un punto o estar vacíos).
\medskip

\textbf{Interpretación:} dibujamos en 2D lugares con la misma altura.
\end{frame}

% 19
\begin{frame}{Pregunta: ¿cómo cambia un mapa de nivel?}
Compara estas dos funciones:
\[
f(x,y)=x^2+y^2,
\qquad
g(x,y)=9-x^2-y^2.
\]
\medskip

En ambas, las curvas de nivel no vacías son circunferencias o un punto.
\medskip

\alert{Pero una forma un valle y la otra una colina. ¿Cuál es cuál?}
\pause
\medskip

En $f$, los valores aumentan al alejarnos del origen; en $g$, disminuyen.
\end{frame}

% 20
\begin{frame}{Aplicación informática: una imagen es una función}
\begin{columns}[c]
\begin{column}{.55\textwidth}
Una imagen en escala de grises asigna una \textbf{intensidad} a cada píxel $(i,j)$.
\medskip

\[
I(i,j)\in\{0,\ldots,255\}.
\]
\medskip

Es una función de dos variables, pero con \textbf{dominio discreto}.
\end{column}
\begin{column}{.43\textwidth}
\pilotoimagen{.90\linewidth}{35mm}{Imagen en grises junto a su cuadrícula ampliada de píxeles}
\end{column}
\end{columns}
\end{frame}

% ===== BLOQUE E. GENERALIZACION Y PRACTICA =====
% 21
\begin{frame}{¿Y si tenemos tres variables de entrada?}
El volumen de una caja depende de su longitud, anchura y altura:
\[
V(x,y,z)=xyz.
\]
\medskip

Ahora la entrada es una terna $(x,y,z)$; la salida sigue siendo \textbf{un número}.
\medskip

\alert{Pregunta:} ¿podemos dibujar la gráfica completa en el espacio tridimensional como antes?
\end{frame}

% 22
\begin{frame}{Funciones de $n$ variables y conjuntos de nivel}
En general, una función real de $n$ variables es
\[
f:D\subseteq\mathbb R^n\longrightarrow\mathbb R.
\]
\medskip

Para $f(x,y,z)=x^2+y^2$, el nivel $k=4$ es
\[
x^2+y^2=4,
\]
\vspace{-1mm}

que describe una \textbf{superficie cilíndrica} en $\mathbb R^3$: $z$ puede valer cualquier número real.
\end{frame}

% 23
\begin{frame}{Nivel 1 · Practica lo esencial}
\textbf{Objetivo:} comprender y aplicar directamente.
\medskip
\begin{enumerate}
\item Para $f(x,y)=2x+y$, calcula $f(1,2)$ y $f(0,3)$.
\medskip
\item Determina el dominio y la imagen de
\[
g(x,y)=\sqrt{4-x^2-y^2}.
\]
\item Describe el nivel $k=1$ de $q(x,y)=x^2+y^2$.
\end{enumerate}
\end{frame}

% 24
\begin{frame}{Nivel 2 · Consolida y representa}
\textbf{Objetivo:} combinar ideas y justificar.
\medskip
\begin{enumerate}
\item Dibuja el dominio de $f(x,y)=\ln(y-x^2)$.
\medskip
\item Halla el dominio de
\[
g(x,y)=\frac{3x-5y}{y-x^2}.
\]
\item Para $h(x,y)=\sqrt{9-x^2-y^2}$, obtén los niveles $k=1$ y $k=2$ e interprétalos.
\end{enumerate}
\end{frame}

% 25
\begin{frame}{Nivel 3 · Desafío (ampliación)}
\textbf{Objetivo:} integrar conocimientos anteriores.
\medskip
\begin{enumerate}
\item Describe los niveles $k=-1,0,1$ de
\[
f(x,y)=x^2-y^2.
\]
\item Encuentra el dominio de
\[
F(x,y,z)=\arcsin(xy)e^{3z}+e^{(y+z)\ln(x+z)}.
\]
\end{enumerate}
\smallskip
\textit{Estos problemas son de ampliación: no son el punto de partida.}
\end{frame}

% 26
\begin{frame}{Cerramos el mapa. ¿Qué viene después?}
Ya podemos:
\begin{itemize}
\item Describir magnitudes que dependen de varias entradas.
\item Determinar dominio, codominio e imagen.
\item Distinguir el dominio de la superficie que dibuja la gráfica.
\item Interpretar las curvas de nivel.
\end{itemize}
\bigskip

\alert{Próxima pregunta:} si caminamos por la montaña, ¿cómo averiguamos la pendiente y la dirección de subida más rápida?
\end{frame}

\end{document}
